\chapter{SQL and Data}\label{ch:iv3:sql-and-data}

\begin{ivtrains}
Writing a query that is right on the awkward rows, and reading one that looks right and is not: joins that lose or multiply rows, the three values of SQL logic, window frames, as-of joins, gaps and bad data. Every table is printed, so you can run the query in your head. A stem names the dialect where the answer depends on it.
\end{ivtrains}

\begin{question}[$\star$ \iqroles{QD, Dev, DS, QR} \iqfirm{any} \iqlevel{junior}]\label{iq:iv3:sql-and-data:1}
Here is a table \texttt{trades}, with \texttt{price} a double and \texttt{qty} an integer. Write a query giving, for each symbol and day, the volume-weighted average price, which is the sum of price times quantity divided by the sum of quantity. Include only symbol-days with more than 2 trades. What does it return here?
\begin{center}\small
\begin{tabular}{llrr}
symbol & day & price & qty\\\hline
AAA & 2026-03-02 & 100.0 & 200\\
AAA & 2026-03-02 & 102.0 & 100\\
AAA & 2026-03-02 & 101.0 & 100\\
BBB & 2026-03-02 & 50.0 & 400\\
BBB & 2026-03-02 & 52.0 & 100\\
AAA & 2026-03-03 & 104.0 & 100
\end{tabular}
\end{center}
\end{question}

\begin{question}[$\star$ \iqroles{QD, Dev, DS, QR} \iqfirm{any} \iqlevel{junior}]\label{iq:iv3:sql-and-data:2}
The left table is \texttt{orders} and the right table is \texttt{fills}. You want every order, with its quantity filled on venue X, or NULL if it has none. What does each query in \cref{lst:iv3:sq:venue} return, and which is right?
\begin{center}\small
\begin{tabular}{rl}
order\_id & symbol\\\hline
1 & AAA\\
2 & BBB\\
3 & CCC
\end{tabular}\qquad
\begin{tabular}{rlr}
order\_id & venue & qty\\\hline
1 & X & 100\\
1 & Y & 50\\
2 & Y & 70
\end{tabular}
\end{center}
\omcode{code/interviews-3/07-sql-and-data/sql/venue_fills.sql}{1}{8}{Two left joins.}\label{lst:iv3:sq:venue}
\end{question}

\begin{question}[$\star$ \iqroles{QD, Dev, DS} \iqfirm{any} \iqlevel{junior}]\label{iq:iv3:sql-and-data:3}
The left table is \texttt{orders} and the right table is \texttt{blocked}. What does \cref{lst:iv3:sq:notin} return, and why? Rewrite it so that it returns every order whose trader is not in the blocked list. Should order 3 be in your answer?
\begin{center}\small
\begin{tabular}{rl}
order\_id & trader\\\hline
1 & ann\\
2 & bob\\
3 & NULL
\end{tabular}\qquad
\begin{tabular}{l}
trader\\\hline
bob\\
NULL
\end{tabular}
\end{center}
\omcode{code/interviews-3/07-sql-and-data/sql/not_in.sql}{1}{3}{Orders from traders not blocked.}\label{lst:iv3:sq:notin}
\end{question}

\begin{question}[$\star$ \iqroles{QD, Dev, DS} \iqfirm{any} \iqlevel{junior}]\label{iq:iv3:sql-and-data:4}
A table \texttt{quotes} has a double column \texttt{bid}, holding the five values 10.0, NULL, 20.0, 20.0, NULL. What do the two statements in \cref{lst:iv3:sq:counts} return, value by value? Which of the last two columns of the first statement is the average bid?
\omcode{code/interviews-3/07-sql-and-data/sql/counts.sql}{1}{4}{Counting with NULLs.}\label{lst:iv3:sq:counts}
\end{question}

\begin{question}[$\star$ \iqroles{QD, Dev, DS} \iqfirm{any} \iqlevel{junior}]\label{iq:iv3:sql-and-data:5}
A table \texttt{fills} has two rows, with integer columns \texttt{filled} and \texttt{ordered} holding (30, 100) and (45, 50). What does \cref{lst:iv3:sq:ratio} return on SQLite, on DuckDB and on PostgreSQL? Write a version that returns the same on all three.
\omcode{code/interviews-3/07-sql-and-data/sql/fill_ratio.sql}{1}{1}{The share of the order filled.}\label{lst:iv3:sq:ratio}
\end{question}

\begin{question}[$\star$ \iqroles{QD, Dev, DS} \iqfirm{any} \iqlevel{junior}]\label{iq:iv3:sql-and-data:6}
A feed sometimes sends a trade again, and sometimes corrects it by sending the same \texttt{trade\_id} with a new quantity. In the table \texttt{trades} below, \texttt{seq} is the arrival order. List the IDs that repeat, and say whether their copies agree. Then give the total quantity per symbol, where the latest version of each trade wins. Would \texttt{SELECT DISTINCT} have given the same totals?
\begin{center}\small
\begin{tabular}{rrlr}
seq & trade\_id & symbol & qty\\\hline
1 & 1 & AAA & 100\\
2 & 2 & AAA & 200\\
3 & 2 & AAA & 200\\
4 & 3 & BBB & 50\\
5 & 3 & BBB & 60\\
6 & 4 & BBB & 70
\end{tabular}
\end{center}
\end{question}

\begin{question}[$\star\star$ \iqroles{QD, Dev, DS, QR} \iqfirm{any} \iqlevel{junior, experienced}]\label{iq:iv3:sql-and-data:7}
Here is a table \texttt{prices}. What does \cref{lst:iv3:sq:bare} do on SQLite, on DuckDB and on PostgreSQL? Write a query for the latest price of each symbol that is correct on all three.
\begin{center}\small
\begin{tabular}{lrr}
symbol & ts & price\\\hline
AAA & 1 & 100.0\\
AAA & 3 & 101.0\\
AAA & 2 & 99.0\\
BBB & 5 & 50.0\\
BBB & 4 & 51.0
\end{tabular}
\end{center}
\omcode{code/interviews-3/07-sql-and-data/sql/latest_bare.sql}{1}{3}{The latest price, perhaps.}\label{lst:iv3:sq:bare}
\end{question}

\begin{question}[$\star\star$ \iqroles{QD, Dev, DS} \iqfirm{any} \iqlevel{junior, experienced}]\label{iq:iv3:sql-and-data:8}
The table \texttt{pnl} has one row per trader. You are asked for the top 2 traders by profit on each desk. Run \texttt{ROW\_NUMBER}, \texttt{RANK} and \texttt{DENSE\_RANK} over \texttt{(PARTITION BY desk ORDER BY pnl DESC)} in your head: what does each give on this table? Which traders does a filter of 2 or less keep under each one? Which would you use, and what must you ask first?
\begin{center}\small
\begin{tabular}{llr}
desk & trader & pnl\\\hline
rates & ann & 50\\
rates & bob & 50\\
rates & cid & 40\\
rates & dan & 30\\
fx & eve & 70\\
fx & fay & 60\\
fx & gus & 60
\end{tabular}
\end{center}
\end{question}

\begin{question}[$\star\star$ \iqroles{QD, Dev, DS, QT} \iqfirm{any} \iqlevel{junior, experienced}]\label{iq:iv3:sql-and-data:9}
The table \texttt{fills} holds signed quantities for one symbol, where \texttt{seq} is the order in which fills occurred, and two fills may share a time \texttt{ts}. A short position is against the rules, so the query in \cref{lst:iv3:sq:position} checks for a negative running position. What does it print, and does it catch the breach? Fix it.
\begin{center}\small
\begin{tabular}{rrr}
seq & ts & qty\\\hline
1 & 10 & 100\\
2 & 20 & $-150$\\
3 & 20 & 100\\
4 & 30 & $-50$
\end{tabular}
\end{center}
\omcode{code/interviews-3/07-sql-and-data/sql/position.sql}{1}{3}{A running position.}\label{lst:iv3:sq:position}
\end{question}

\begin{question}[$\star\star$ \iqroles{QD, Dev, DS} \iqfirm{any} \iqlevel{junior, experienced}]\label{iq:iv3:sql-and-data:10}
The left table is \texttt{orders} and the right table is \texttt{fills}. \Cref{lst:iv3:sq:fanout} is meant to give, for each desk, the total quantity ordered and the total filled. What does it print, what is wrong with each number, and how do you fix it?
\begin{center}\small
\begin{tabular}{rlr}
order\_id & desk & qty\\\hline
1 & rates & 100\\
2 & rates & 50\\
3 & fx & 200
\end{tabular}\qquad
\begin{tabular}{rr}
order\_id & filled\\\hline
1 & 60\\
1 & 40\\
3 & 150
\end{tabular}
\end{center}
\omcode{code/interviews-3/07-sql-and-data/sql/fanout.sql}{1}{4}{Ordered and filled, by desk.}\label{lst:iv3:sq:fanout}
\end{question}

\begin{question}[$\star\star$ \iqroles{QD, Dev, DS, QR} \iqfirm{any} \iqlevel{junior, experienced}]\label{iq:iv3:sql-and-data:11}
In the table \texttt{closes}, \texttt{day} is a calendar day number, and days 4 and 5 are a weekend with no row. What are the two columns of averages in \cref{lst:iv3:sq:moving}, row by row? Which one is a three-day moving average?
\begin{center}\small
\begin{tabular}{rr}
day & price\\\hline
1 & 10.0\\
2 & 20.0\\
3 & 30.0\\
6 & 60.0\\
7 & 90.0
\end{tabular}
\end{center}
\omcode{code/interviews-3/07-sql-and-data/sql/moving.sql}{1}{5}{Two moving averages.}\label{lst:iv3:sq:moving}
\end{question}

\begin{question}[$\star\star$ \iqroles{QD, Dev, DS, QR} \iqfirm{any} \iqlevel{junior, experienced}]\label{iq:iv3:sql-and-data:12}
\Cref{lst:iv3:sq:barsbad} should turn the table \texttt{ticks} into one bar per minute, with open, high, low and close. How many rows does it return, and what are they? Why? Fix it.
\begin{center}\small
\begin{tabular}{rrr}
minute & seq & price\\\hline
1 & 1 & 100.0\\
1 & 2 & 103.0\\
1 & 3 & 101.0\\
2 & 4 & 102.0\\
2 & 5 & 99.0
\end{tabular}
\end{center}
\omcode{code/interviews-3/07-sql-and-data/sql/bars_bad.sql}{1}{8}{Bars from ticks.}\label{lst:iv3:sq:barsbad}
\end{question}

\begin{question}[$\star\star$ \iqroles{QD, Dev, DS, QR} \iqfirm{any} \iqlevel{junior, experienced}]\label{iq:iv3:sql-and-data:13}
The left table is \texttt{trades} and the right table is \texttt{quotes}, with \texttt{ts} in seconds. For each trade, return the bid and ask of the latest quote at or before the trade, but treat a quote more than 10 seconds old as missing. Use a correlated subquery, not any special join syntax. What do you return for each trade here? What does your query cost with and without an index on \texttt{quotes.ts}?
\begin{center}\small
\begin{tabular}{rr}
trade\_id & ts\\\hline
1 & 2\\
2 & 5\\
3 & 12\\
4 & 30
\end{tabular}\qquad
\begin{tabular}{rrr}
ts & bid & ask\\\hline
3 & 99.9 & 100.1\\
10 & 100.3 & 100.6\\
12 & 100.4 & 100.7
\end{tabular}
\end{center}
\end{question}

\begin{question}[$\star\star$ \iqroles{QD, Dev, LLE, DS} \iqfirm{market maker, crypto, hedge fund} \iqlevel{junior, experienced}]\label{iq:iv3:sql-and-data:14}
A table \texttt{msgs} has one integer column, \texttt{seq}, the sequence numbers received from a feed that should have no holes. Its rows are 1, 2, 3, 6, 7 and 11. Write a query that returns each missing range as a first and a last number. A shorter query, the largest minus the smallest plus one minus the row count, gives the number missing: why is that not enough? What can no query on this table tell you?
\end{question}

\begin{question}[$\star\star$ \iqroles{QD, Dev, DS, QR} \iqfirm{market maker, crypto, hedge fund} \iqlevel{junior, experienced}]\label{iq:iv3:sql-and-data:15}
The table \texttt{quotes} below holds one symbol, in arrival order \texttt{seq}. Write one query that returns each suspect row and one reason for it, checking in this order: bid at or above ask; a size of zero or less; a time that does not increase; and a mid-price that moved more than 5\% from the previous row's. Which rows come back? Which of them is a good row? How would you stop one bad print from flagging two rows?
\begin{center}\small
\begin{tabular}{rrrrr}
seq & ts & bid & ask & bid\_size\\\hline
1 & 1 & 100.0 & 100.2 & 500\\
2 & 2 & 100.1 & 100.3 & 400\\
3 & 3 & 100.4 & 100.2 & 300\\
4 & 4 & 100.1 & 100.3 & 0\\
5 & 4 & 100.1 & 100.3 & 200\\
6 & 6 & 110.0 & 110.2 & 100\\
7 & 7 & 100.2 & 100.4 & 300
\end{tabular}
\end{center}
\end{question}

\begin{question}[$\star\star$ \iqroles{QD, Dev, LLE, DS} \iqfirm{any} \iqlevel{junior, experienced}]\label{iq:iv3:sql-and-data:16}
A table \texttt{latency} has one integer column \texttt{us}, holding 10, 12, 13, 15, 20, 22, 25, 30, 80 and 200. On DuckDB, what do \texttt{quantile\_disc(us, 0.9)} and \texttt{quantile\_cont(us, 0.9)} return, and why do they differ? SQLite 3.37 has neither, so write the 90th percentile there by the nearest-rank rule: the smallest value with at least 90\% of the rows at or below it. Which definition would you report for latency?
\end{question}

\begin{question}[$\star\star$ \iqroles{QD, Dev, DS} \iqfirm{any} \iqlevel{junior, experienced}]\label{iq:iv3:sql-and-data:17}
On SQLite, a table \texttt{trades} has the index in \cref{lst:iv3:sq:plans}, on \texttt{(symbol, ts)}. For each of the five queries, does the plan search the index or scan the table, and does it need a separate sort? Explain the rule behind your answers.
\omcode{code/interviews-3/07-sql-and-data/sql/index_plans.sql}{1}{13}{One index, five queries.}\label{lst:iv3:sq:plans}
\end{question}

\begin{question}[$\star\star\star$ \iqroles{QD, Dev, DS, QR} \iqfirm{any} \iqlevel{junior, experienced}]\label{iq:iv3:sql-and-data:18}
The table \texttt{daily} holds one row per trading day, in which \texttt{day} increases but need not be consecutive. Write a query that returns the first day, the last day and the length of the longest run of consecutive rows with positive profit. On ties, take the earliest. What does it return here? Use no loop or recursion.
\begin{center}\small
\begin{tabular}{rr}
day & pnl\\\hline
1 & 5\\
2 & 3\\
3 & $-2$\\
6 & 1\\
7 & 2\\
8 & 4\\
9 & $-1$\\
10 & 7
\end{tabular}
\end{center}
\end{question}

\begin{question}[$\star\star\star$ \iqroles{QD, Dev, DS} \iqfirm{any} \iqlevel{junior, experienced}]\label{iq:iv3:sql-and-data:19}
The table \texttt{events} holds a trader's actions, with \texttt{ts} in seconds. A session is a run of one trader's events in which each follows the last by 300 seconds or less, and a longer gap starts a new session. Write a query that returns, for each session, the trader, the session number, the first and last times and the number of events. What does it return here? Why is a gap of exactly 300 the case to test?
\begin{center}\small
\begin{tabular}{lr}
trader & ts\\\hline
ann & 100\\
ann & 200\\
ann & 700\\
ann & 800\\
ann & 1500\\
bob & 150\\
bob & 460\\
bob & 760
\end{tabular}
\end{center}
\end{question}

\begin{question}[$\star\star\star$ \iqroles{QD, Dev, DS, QR} \iqfirm{any} \iqlevel{junior, experienced}]\label{iq:iv3:sql-and-data:20}
In the table \texttt{closes}, a NULL price means no trade that day. Fill each NULL with the last known price, and leave a NULL with nothing before it as it is. Write it for SQLite 3.37, which does not support \texttt{IGNORE NULLS}, and say how DuckDB can do it in one line. What do you return here? Why would you not fill backwards as well?
\begin{center}\small
\begin{tabular}{rr}
day & price\\\hline
1 & NULL\\
2 & 10.0\\
3 & NULL\\
4 & NULL\\
5 & 12.0\\
6 & NULL
\end{tabular}
\end{center}
\end{question}

\begin{question}[$\star\star\star$ \iqroles{QD, Dev, DS, QR} \iqfirm{market maker, hedge fund, crypto} \iqlevel{experienced}]\label{iq:iv3:sql-and-data:21}
The left table is \texttt{trades} and the right table is \texttt{quotes}, each with millions of rows per symbol, and neither has an index. For each trade, return the \texttt{mid} of the latest quote of the same symbol at or before the trade, or NULL if none. A quote at the same time as a trade counts. Give the query on DuckDB, using its special join. Then give a query with no such join and no correlated subquery, which runs on SQLite 3.37. What does each cost, and what do you return here?
\begin{center}\small
\begin{tabular}{lr}
symbol & ts\\\hline
AAA & 5\\
AAA & 12\\
BBB & 1\\
BBB & 7
\end{tabular}\qquad
\begin{tabular}{lrr}
symbol & ts & mid\\\hline
AAA & 3 & 100.0\\
AAA & 12 & 101.0\\
BBB & 2 & 50.0\\
BBB & 9 & 51.0
\end{tabular}
\end{center}
\end{question}

\begin{question}[$\star\star\star$ \iqroles{QD, Dev, DS, QR} \iqfirm{hedge fund, bank} \iqlevel{experienced}]\label{iq:iv3:sql-and-data:22}
A ticker can pass from one company to another. The left table, \texttt{master}, gives each instrument's ticker and the dates it held it, from \texttt{valid\_from} up to but not including \texttt{valid\_to}, and the right table is \texttt{trades}. Dates are ISO text. \Cref{lst:iv3:sq:pit} is a join someone wrote. What does it return, and what are the two separate problems with this result? Fix the query, and write a data-quality query that finds the bad rows in \texttt{master}.
\begin{center}\small
\begin{tabular}{rlll}
inst\_id & ticker & valid\_from & valid\_to\\\hline
1 & ABC & 2020-01-01 & 2024-05-01\\
2 & ABC & 2024-05-01 & 9999-12-31\\
3 & XYZ & 2020-01-01 & 2023-01-01\\
4 & XYZ & 2022-06-01 & 9999-12-31
\end{tabular}\qquad
\begin{tabular}{llr}
ticker & day & qty\\\hline
ABC & 2024-04-30 & 10\\
ABC & 2024-05-01 & 20\\
XYZ & 2022-07-01 & 30
\end{tabular}
\end{center}
\omcode{code/interviews-3/07-sql-and-data/sql/pit_between.sql}{1}{4}{A join on validity dates.}\label{lst:iv3:sq:pit}
\end{question}

\begin{question}[$\star\star\star$ \iqroles{QD, Dev, DS, QR} \iqfirm{market maker, hedge fund, crypto} \iqlevel{experienced}]\label{iq:iv3:sql-and-data:23}
The table \texttt{trades} has an integer time \texttt{ts} in seconds, a double \texttt{price} and an integer \texttt{qty}. For each trade, give the volume-weighted average price over the five minutes ending at that trade, that is, over trades with $t - 300 < \text{ts} \le t$. Use a window function with a frame in time, not in rows. What do you return here? Two trades share the time 420: do they see each other? What must change if \texttt{ts} were a double?
\begin{center}\small
\begin{tabular}{rrr}
ts & price & qty\\\hline
0 & 100.0 & 10\\
120 & 101.0 & 10\\
299 & 102.0 & 20\\
300 & 103.0 & 10\\
420 & 100.0 & 10\\
420 & 99.0 & 30
\end{tabular}
\end{center}
\end{question}

\begin{question}[$\star\star\star\star$ \iqroles{QD, Dev, DS, QT} \iqfirm{any} \iqlevel{experienced}]\label{iq:iv3:sql-and-data:24}
The table \texttt{fills} holds one symbol's buys and sells, in order \texttt{seq}, and the position never goes negative. Sales are matched to buys first in, first out: each sale uses up the oldest shares still held. In one SQL statement, with window functions but no loop and no recursion, return one row per matched pair of buy and sell, with the quantity and the profit, on SQLite 3.37 and on DuckDB. What is the total profit here? Does your query still hold if the position can go short?
\begin{center}\small
\begin{tabular}{rlrr}
seq & side & price & qty\\\hline
1 & B & 100 & 100\\
2 & B & 110 & 50\\
3 & S & 120 & 120\\
4 & B & 90 & 30\\
5 & S & 130 & 60
\end{tabular}
\end{center}
\end{question}

\begin{question}[$\star\star\star\star$ \iqroles{QD, Dev, DS} \iqfirm{market maker, crypto} \iqlevel{experienced}]\label{iq:iv3:sql-and-data:25}
The table \texttt{events} is the history of the buy side of one book: at step \texttt{seq}, the resting quantity at \texttt{price} changes by \texttt{delta}. In one SQL statement, on SQLite 3.37 and on DuckDB, return for each step the best bid after that step, meaning the highest price with a positive quantity, and the quantity there, or NULL if the side is empty. What do you return here? What does your query cost, and would you run it this way on a day of data?
\begin{center}\small
\begin{tabular}{rrr}
seq & price & delta\\\hline
1 & 100 & 50\\
2 & 101 & 20\\
3 & 100 & 30\\
4 & 101 & $-20$\\
5 & 99 & 40\\
6 & 100 & $-80$
\end{tabular}
\end{center}
\end{question}
